\documentclass[11pt,a4paper]{article}
\usepackage[margin=23mm]{geometry}
\usepackage[T1]{fontenc}
\usepackage{lmodern,xurl,hyperref}
\hypersetup{hidelinks}
\setlength{\parindent}{0pt}\setlength{\parskip}{7pt}
\emergencystretch=2em
\begin{document}
\textbf{Revised manuscript submission to Scientific Reports}\\
Submission ID: \texttt{20e66410-c472-42dc-ab99-957c24a43f96}\\
To: Dr Osman Orhan, Editorial Board Member

Dear Dr Orhan,

We present our revised manuscript, \emph{Spatial aggregation alters population-weighted estimates of InSAR observation support}, for reconsideration. The original title was \emph{Auditing observability bias in InSAR-derived subsidence exposure estimates with a reproducible geospatial workflow}.

The revised study quantifies how spatial aggregation changes population-weighted InSAR support summaries when the support field, domain and population total are fixed. Within-cell covariance explains the aggregation effect; controlled missingness identifies both error directions and cases with small aggregate bias. Crossing GHSL and WorldPop with 28- and 224-pair selections preserves a positive difference across all 16 tested cases while demonstrating strong population-model dependence.

A connected LiCSBAS experiment distinguishes sampled input support from final processing quality and evaluates internal reconstruction using excluded edges. The Fresno analyses test transfer to published observation polygons and assess population-model differences against Census and ACS reporting units. Further analyses quantify estimator availability under clustered missingness, the coverage--residual tradeoff and regional limits to transfer. The smaller regional effects, incomplete residential coverage and mixed GNSS agreement limit the scope of the conclusions. We corrected estimator-availability accounting, DWR interval annualization, population denominators and the interpretation of the block-group benchmark.

We have withdrawn the interpretation of the original 3.77-million figure as a verified count of people exposed to subsidence in unobserved locations. The revised conclusion concerns allocation and aggregation effects; it does not establish deformation in missing locations. The Iran source attribution has also been corrected, and its incomplete analysis is excluded from quantitative applications.

The accompanying files contain the clean editable manuscript, Supplementary Information, separate main figures and a point-by-point response with final section and PDF-page locations. The optional marked manuscript shows changes from the preceding local 1 October 2026 writing package. The scientific reproduction package and additional official source inputs are publicly available at the fixed GitHub release \href{https://github.com/niubisile-wq/openinsar-observability-bias/releases/tag/v0.6.0}{v0.6.0}. The previous Zenodo DOI covers v0.3.0 only; a revised DOI deposit remains pending.

Sincerely,\\
Zixuan Liu and Wei Xiong\\
Corresponding author: Wei Xiong, \href{mailto:xw@mail.hbut.edu.cn}{xw@mail.hbut.edu.cn}\\
School of Electrical and Electronic Engineering, Hubei University of Technology, Wuhan 430068, China; Department of Computer Science and Engineering, University of South Carolina, Columbia, SC 29201, USA.
\end{document}
